\subsection{The Culture Conversion Model}
\label{subsec:cult}

The cultural capital approach to taste formation provides a very different picture of the stability of tastes than that which can be gleaned from contemporary network theory.  Instead of implying that tastes are unstable, ``shallow'' contents liable to be replaced at the same pace that network relations are replaced, the cultural capital approach suggest that taste are ``deep'' durable dispositions and trans-temporally stable forms of embodied cultural capital \citep{bourdieu86, holt98}.  It is the relative stability of tastes vis a vis network relations that allows them to function as resources to help sustain and maintain these network relations over time---primarily by serving as topical resources in conversational interaction rituals \citep{dimaggio87, lizardo06a} providing a common focus of attention and affective commitment \citep{collins04}.  This helps to shift interactions away from purely ``instrumental'' purposes of information gathering and rational goal attainment to sociable interaction for its own sake \citep{granovetter02, simmel49}. 

\citet{lizardo06a} uses this stability assumption to suggest that tastes can in fact be seen as probable (reciprocal) drivers of network connections rather than being seen as exclusively determined by those connections.  Using two-stage instrumental variables and matching methods on cross-sectional data, that study uncovers evidence consistent with this reciprocal effect. Engagement in a wide variety of cultural practices increases network size net of standard controls and after statistically correcting for simultaneity bias.  This is in agreement with dynamic theories of cultural and network evolution, which rather than postulating a one-way avenue of causation going from networks to cultural tastes propose a reciprocal (over-time) two-way casual interaction between the two \citep{carley91, mark98a}. However, in the aforementioned study the stability hypothesis remains an empirically unverified background assumption.  Furthermore to reliance on cross-sectional data, the casual effect of cultural taste on network maintenance or even the more widely accepted effect of network turnover on cultural taste loss, has not been decisively established.
 
Thus, according to the cultural-capital theory inspired  conversion model of cultural into social capital \citep{dimaggio_mohr85, dimaggio_useem78a, lizardo06a}, the primary ways in which tastes are used to sustain social networks, is by allowing the individual to connect to alters that are both relationally close---when specialist or ``niche'' cultural tastes are concerned---and relationally distant---when talking about aptitude for ``popular'' cultural forms.  From this perspective, culture consumption and the ability to decode and appropriate symbolic goods produced in the urban, folk and ``culture-industry'' realms \citep{crane93} are important resources that individuals mobilize in the construction and maintenance of active social relationships \citep{cardon_granjon05, dimaggio87}.  The culture conversion model can in this particular respect be thought of as an elaboration and extension (to the realm of cultural taste) of  the classical \citep[e.g.]{lazarsfeld_merton54} process model of friendship formation, with similarity based on cultural taste and practices \citep{dimaggio87, mark03} playing the role that ``value-homophily'' played in the original formulation.    In this respect,  network models of taste formation through social network ties \citep{mark98b} can be thought of as partially based on the Lazarsfeld-Merton formulation.  However, these models raise what was explicitly a secondary dynamic in the original scheme---i.e. mutual ``value-adjustment'' through social interaction of initially different value dispositions \citep[33]{lazarsfeld_merton54}---to being the only operative, and thus primary dynamic.  It is clear however that for \citet[22--29]{lazarsfeld_merton54} ``selective'' processes whereby contacts are sifted based on similarity of attitudinal dispositions similar to those emphasized in the culture-conversion account were considered primary.

There is plenty of evidence that shows that tastes do play the role accorded to them (vis a vis network ties) in the culture conversion model.  As \citet[302]{cardon_granjon05}  note in their recent study of the role of culture consumption in the formation of social connections among young adults in France, ``{\ldots}a large proportion of ordinary conversational interaction is generally devoted to comments{\ldots}on cultural practices and recreational consumption. These are recurrent topics, present in a wide range of social situations{\ldots}Many social relations are constructed, in different ways, around cultural and recreational activities.''  As \citep[453]{fine79} noted in a classic paper, ``culture produces considerable discussion, and thus {\ldots}[its] content {\ldots}will become a resource in conversational interaction.'' Accordingly, those cultural goods which find an audience, will be deployed as the basis of conversation and will read the ``fodder for least-common denominator talk.'' We should expect that the consumption of aesthetic goods infuses local conversation-based interaction rituals with a common focus of attention allowing persons to spend time engaged in (Simmelian) interaction for its own sake\citeyearpar[443]{dimaggio87}. 

For instance, Elizabeth Long, in her study of women's reading groups in Houston, Texas, finds that cultural knowledge and taste become the primary content of social interaction. Consistent with the claim that connects the consumption of aesthetic products and sociability, she finds that reading group members ``tend to press books into service for the meanings that they transmit and the conversations they generate'' \citeyearpar[73]{long03}.  The primary function of aesthetic consumption experiences is to serve as the non-instrumental content that justifies the existence of status-segregated conversational rituals: ``{\ldots}talk about entertainment-the discussions of movies, music, sports, actors and athletes…make up so much of casual chitchat.'' This becomes even more important ``{\ldots}if we include talk about things people consume or seek out-food, clothes, houses and gardens, restaurants scenery, vacations. In some social circles, this entertainment talk makes up most of the content of conversations; in all classes, it is important" \citeyearpar[175]{collins75}. The Simmelian language of “social circles” is not accidental, since as \cite[797]{kadushin66} noted in his classic study on the subject, Simmelian urban social circles are produced and reproduced through shared cultural tastes---e.g. concert going, visiting the museum, etc.---and other activities intrinsically connected to culture-mediated interaction. 

From this perspective of the culture-conversion model therefore, while the effect of social network change on cultural taste loss is not denied, whatever relational (or “network”) effects on taste exist, must operate in the background of high-levels of cross-temporal taste stability and the reciprocal action of cultural taste on social networks.

 